% !TeX program = xelatex
% Сборка из папки mathan: xelatex report_template.tex (два раза).
% Рисунки и таблицы: python examples/report_example.py
% Для собственного отчёта сохранить копию под именем report.tex.
\documentclass[12pt,a4paper]{article}
\usepackage{fontspec}
\usepackage{polyglossia}
\setdefaultlanguage{russian}
\setmainfont{Times New Roman}
\setmonofont{Courier New}
\newfontfamily\cyrillicfonttt[Script=Default]{Courier New}
\usepackage{amsmath,amssymb,graphicx,array}
\usepackage[margin=2cm]{geometry}
\usepackage{xcolor,listings}
\usepackage[unicode,hidelinks]{hyperref}
\setlength{\parindent}{0pt}
\setlength{\parskip}{6pt}
\setlength{\emergencystretch}{3em}
\renewcommand{\arraystretch}{1.25}
\newcommand{\taskhead}[1]{\par\textbf{#1}\par}
% Сохранять ASCII-дефис U+002D при копировании кода из PDF.
\lstset{literate={-}{{\char45}}1}
\lstdefinestyle{commands}{
  basicstyle=\ttfamily\small,
  backgroundcolor=\color{gray!10},
  frame=single,framerule=0pt,framesep=6pt,
  xleftmargin=6pt,xrightmargin=6pt,
  breaklines=true,columns=fullflexible,keepspaces=true,
  showstringspaces=false
}

\begin{document}
\sloppy
\begin{titlepage}
{\color{black!50}
\textbf{\textit{Примечание:}} В данном документе представлен сокращённый
образец оформления титульного листа и результатов выполнения лабораторной
работы для задачи 2.2 и одного эксперимента из подраздела 3.5
(поиск минимума затрат градиентным спуском).
Приведены примеры аналитических вычислений, вызовов функций,
таблиц сравнения и проверок, рисунков и выводов.
Сведения об участниках приведены для примера; в вычислениях использованы
параметры варианта 1 из приложений А и Б условия.
\par}
\vspace{1.5cm}
{\large\bfseries Отчёт по лабораторной работе\par}
\textbf{По дисциплине:} Математический анализ (3 семестр)

\textbf{На тему:} Численное вычисление кратных интегралов.
Исследование и оптимизация динамической системы
\vspace{1cm}

\textbf{Вариант № 1}

\textbf{Выполнили} студенты гр. № 123456:

\textit{Иванов И.И., Петров П.П., Сидоров С.С.}

\textbf{Дата выполнения:} [дата]
\vspace{0.7cm}

\textbf{Цель работы}

Изучение роли замены координат в вычислении тройного интеграла,
точности численного интегрирования и влияния шага градиентного спуска
на поиск оптимального распределения ресурса.
\vspace{0.5cm}

\textbf{Условия выполнения}

ОС: macOS (Darwin 25.2.0). Python 3.12.3. NumPy 2.2.6.
Matplotlib 3.10.6. Сборка отчёта: XeLaTeX.
Результаты получены запуском демонстрационного скрипта
\texttt{examples/report\_example.py}. В собственном отчёте указываются
фактическое окружение и результаты собственной реализации.
\end{titlepage}

\taskhead{Порядок выполнения}
\begin{enumerate}
\item[2.2] Область описана в цилиндрических координатах, вычислен якобиан.
Интеграл вычислен аналитически и методом средних прямоугольников.
Построены изображение области и график ошибки; результаты сопоставлены
в таблице. Выполнена проверка объёма.
\item[3.5] Для задачи оптимизации затрат вычислено аналитическое решение.
Выполнен градиентный спуск с постоянным и убывающим шагами.
Построены траектории на фоне линий уровня, сопоставлены точки минимума,
затраты и числа итераций.
\end{enumerate}

\newpage
\setcounter{section}{1}
\section{Тройные интегралы и замена координат}
\setcounter{subsection}{1}
\subsection{Цилиндрические координаты}
\taskhead{Исходные данные и аналитическое решение}
Для варианта 1 заданы $R=1$, $H=2$. Рассматривается интеграл
\[
I=\iiint_\Omega(1+x^2+y^2+z)\,dx\,dy\,dz,
\quad \Omega=\{x^2+y^2\leq1,\ y\geq0,\ 0\leq z\leq2\}.
\]
Область представляет половину цилиндра радиуса 1 и высоты 2.
В цилиндрических координатах
\[
x=r\cos\theta,\quad y=r\sin\theta,\quad z=z,
\qquad 0\leq r\leq1,\quad0\leq\theta\leq\pi,\quad0\leq z\leq2.
\]
Условие $y\geq0$ ограничивает угол полуоборотом. Матрица частных
производных даёт якобиан
\[
J=\det\begin{pmatrix}
\cos\theta&-r\sin\theta&0\\
\sin\theta&r\cos\theta&0\\0&0&1
\end{pmatrix}=r(\cos^2\theta+\sin^2\theta)=r.
\]
Следовательно, $dx\,dy\,dz=r\,dr\,d\theta\,dz$, а преобразованная
функция с учётом якобиана равна $g(r,\theta,z)=(1+r^2+z)r$.
Последовательное интегрирование даёт
\[
\begin{aligned}
I&=\int_0^\pi\int_0^1\int_0^2(1+r^2+z)r\,dz\,dr\,d\theta\\
&=\pi\int_0^1\left[(1+r^2)z+\frac{z^2}{2}\right]_0^2r\,dr\\
&=\pi\int_0^1(4r+2r^3)\,dr
=\pi\left[2r^2+\frac{r^4}{2}\right]_0^1
=\frac{5\pi}{2}\approx7.853981634.
\end{aligned}
\]
Для проверки объёма исходная функция заменяется на 1, а якобиан сохраняется:
\[
V=\int_0^\pi\int_0^1\int_0^2r\,dz\,dr\,d\theta
=\pi=\frac{\pi R^2H}{2}.
\]

\taskhead{Программный результат и сравнение}
Использована трёхмерная формула средних прямоугольников.
При $n$ отрезках по каждой координате шаги равны
$h_r=1/n$, $h_\theta=\pi/n$, $h_z=2/n$; значения $g$ вычисляются
в центрах всех $n^3$ ячеек.
Вызов демонстрационной функции из папки лабораторной:
\begin{lstlisting}[style=commands]
from examples.report_example import midpoint_cylinder

print(midpoint_cylinder(8))               # 7.841709787671397
print(midpoint_cylinder(8, volume=True))  # 3.141592653589793
\end{lstlisting}
В собственном отчёте приводится вызов реализованной студентом функции
из \texttt{solution.py}. Все значения образца воспроизводятся командой
\begin{lstlisting}[style=commands]
python examples/report_example.py
\end{lstlisting}
Для результата $I_n$ абсолютная ошибка равна $\Delta_n=|I_n-I|$,
относительная --- $\delta_n=\Delta_n/|I|$.
Например, при $n=8$:
\[
\Delta_8=\left|7.841709787671397-\frac{5\pi}{2}\right|
\approx0.0122718463,\qquad
\delta_8\approx\frac{0.0122718463}{7.853981634}=0.0015625.
\]
Ошибки рассчитаны до округления. Результаты приведены
в таблице~\ref{tab:integrals}; $h=1/n$ --- нормированный шаг.
Запись \texttt{e-03} в таблицах означает множитель $10^{-3}$.
\begin{table}[htbp]
\centering\small
\caption{Сравнение численного интеграла с $I=5\pi/2$}
\label{tab:integrals}
\begin{tabular}{|r|r|r|r|r|}\hline
$n$ & $h$ & $I_n$ & $\Delta_n$ & $\delta_n$\\\hline
8 & 0.125000 & 7.841709788 & 1.227185e-02 & 1.562500e-03 \\
16 & 0.062500 & 7.850913672 & 3.067962e-03 & 3.906250e-04 \\
32 & 0.031250 & 7.853214644 & 7.669904e-04 & 9.765625e-05 \\
64 & 0.015625 & 7.853789886 & 1.917476e-04 & 2.441406e-05 \\
\hline\end{tabular}
\end{table}
Объём на всех четырёх сетках получен равным $3.141592653589793$:
расхождение с числовым значением $\pi$ в выполненном расчёте равно нулю.
Это согласуется с тем, что преобразованная функция объёма $r$ линейна
и метод средних прямоугольников интегрирует её точно с учётом округления.

\clearpage
\taskhead{Графическая проверка}
На рисунке~\ref{fig:cylinder} показана область интегрирования с равным
масштабом пространственных осей. На рисунке~\ref{fig:error}
представлена зависимость абсолютной ошибки от шага в логарифмическом масштабе.
\begin{figure}[htbp]
\centering
\includegraphics[width=0.67\linewidth]{examples/report_figures/cylinder.pdf}
\caption{Половина цилиндра при $R=1$, $H=2$}
\label{fig:cylinder}
\end{figure}
\begin{figure}[htbp]
\centering
\includegraphics[width=0.75\linewidth]{examples/report_figures/integration_error.pdf}
\caption{Абсолютная ошибка метода средних прямоугольников}
\label{fig:error}
\end{figure}
\taskhead{Вывод}
При удвоении числа разбиений ошибка уменьшается в четыре раза,
что согласуется со вторым порядком метода. При $n=64$ относительная
ошибка составляет примерно $0.0024414\%$. Объём совпал с геометрическим
значением; это подтверждает корректность границ и множителя якобиана
в данной проверке, но само по себе не доказывает точность основного интеграла.

\clearpage
\section{Исследование и оптимизация динамической системы}
\setcounter{subsection}{4}
\subsection{Численная проверка}
\taskhead{Исходные данные и аналитическое решение}
Рассмотрен только пункт 5: градиентный спуск для двух звеньев без ограничения
на сумму ресурсов. Для варианта 1:
\[
\lambda=400,\quad\mu_1=50,\quad\mu_2=80,\quad a=1,\quad b=2.
\]
Целевая функция и допустимая область имеют вид
\[
\Phi(c_1,c_2)=c_1+c_2+\frac{800}{50c_1-400}
+\frac{800}{80c_2-400},\qquad c_1>8,\ c_2>5.
\]
Частная производная по ресурсу звена $i$ равна
\[
\Phi_{c_i}=a-\frac{b\lambda\mu_i}{(\mu_ic_i-\lambda)^2}.
\]
Из равенства производной нулю и положительности запаса мощности следует
\[
c_i^*=\frac\lambda{\mu_i}+\sqrt{\frac{b\lambda}{a\mu_i}}.
\]
Подстановка параметров даёт
\[
c_1^*=8+\sqrt{16}=12,\qquad c_2^*=5+\sqrt{10}\approx8.162277660.
\]
Вторые производные $\Phi_{c_ic_i}=2b\lambda\mu_i^2/(\mu_ic_i-\lambda)^3$
положительны, смешанные равны нулю. Функция строго выпукла в допустимой
области, поэтому найденная точка является единственным глобальным минимумом.
Подстановка её координат в функцию даёт
\[
\Phi^*=21+2\sqrt{10}\approx27.324555320.
\]

\taskhead{Параметры численного эксперимента}
Начальная точка $c^{(0)}=(10,10)$ допустима и имеет затраты
$\Phi(c^{(0)})=20+800/100+800/400=30$.
Градиент оценивался центральными разностями с шагом $h_d=10^{-5}$
по каждой координате. Здесь $h_d$ --- шаг дифференцирования,
а не шаг оптимизации. Сравнивались два правила:
\[
\eta_j=0.2,\qquad \eta_j=\frac{0.4}{1+0.01j},\qquad j=0,1,\ldots.
\]
На каждой итерации использовалось обновление
$c^{(j+1)}=c^{(j)}-\eta_j\widehat{\operatorname{grad}}\Phi(c^{(j)})$,
где знак над градиентом означает численную оценку.
Остановка выполнялась при максимальном модуле компоненты численного
градиента не более $10^{-7}$; предел --- 10000 обновлений.
Допустимость проверялась на каждом шаге; в обоих запусках нарушений не было.
Число итераций ниже означает число выполненных обновлений точки.

Например, точный градиент в начальной точке равен
\[
\operatorname{grad}\Phi(10,10)
=\left(1-\frac{800\cdot50}{100^2},\
1-\frac{800\cdot80}{400^2}\right)=(-3,0.6).
\]
При постоянном шаге первое обновление даёт приближённо
$(10,10)-0.2(-3,0.6)=(10.6,9.88)$.

\taskhead{Программный результат и сравнение}
Вызов демонстрационной реализации:
\begin{lstlisting}[style=commands]
from examples.report_example import optimize, cost

point, path = optimize("constant")
print(point)
print(cost(point))
print(len(path) - 1)
\end{lstlisting}
Численные результаты представлены в таблице~\ref{tab:optimization}.
Для сравнения координат использованы
\[
\Delta_c=\max_{i=1,2}|\widehat c_i-c_i^*|,\qquad
\delta_c=\frac{\Delta_c}{\max_{i=1,2}|c_i^*|}=\frac{\Delta_c}{12},
\]
где $\widehat c_i$ --- вычисленные координаты, $c_i^*$ --- аналитические.
\begin{table}[htbp]
\centering\small
\caption{Результаты градиентного спуска}
\label{tab:optimization}
\begin{tabular}{|l|r|r|r|r|r|}\hline
Шаг & Итерации & $\widehat c_1$ & $\widehat c_2$ & $\Phi(\widehat c)$ & $\Delta_c$\\\hline
Постоянный & 146 & 11.999999818 & 8.162277674 & 27.324555320 & 1.824e-07 \\
Убывающий & 101 & 11.999999817 & 8.162277678 & 27.324555320 & 1.829e-07 \\
\hline\end{tabular}
\end{table}
Для постоянного шага $\delta_c\approx1.5203\cdot10^{-8}$,
для убывающего --- $1.5239\cdot10^{-8}$.
Абсолютные ошибки затрат в выполненных расчётах меньше $2\cdot10^{-14}$,
относительные --- меньше $8\cdot10^{-16}$; это уровень машинного округления.
Максимальные модули компонент численного градиента составили
$9.13\cdot10^{-8}$ и $9.17\cdot10^{-8}$ соответственно.
Обе точки удовлетворяют выбранному критерию остановки.

\begin{figure}[htbp]
\centering
\includegraphics[width=0.85\linewidth]{examples/report_figures/gradient_paths.pdf}
\caption{Траектории спуска на фоне линий уровня затрат}
\label{fig:paths}
\end{figure}
На рисунке~\ref{fig:paths} отмечены начальная точка и аналитический минимум;
числа на линиях уровня обозначают значения функции $\Phi$.
Обе траектории приближаются к одному минимуму, оставаясь в допустимой области.
\taskhead{Вывод}
В обоих экспериментах максимальная ошибка координат меньше $2\cdot10^{-7}$.
Убывающий шаг потребовал 101 обновление, постоянный --- 146.
Это сравнение относится к указанным правилам шага и начальной точке;
оно не означает, что убывающий шаг всегда быстрее постоянного.

\taskhead{Проверка реализации}
Таблица~\ref{tab:checks} содержит проверки, выполненные демонстрационным
скриптом. Команда запуска для всех строк:
\begin{lstlisting}[style=commands]
python examples/report_example.py
\end{lstlisting}
Готовый набор автоматических тестов условием не предусмотрен.
Приведённые проверки являются частью данного демонстрационного эксперимента.
\begin{table}[htbp]
\centering\small
\caption{Результаты проверок}
\label{tab:checks}
\begin{tabular}{|p{2.1cm}|p{6.3cm}|p{5.4cm}|}\hline
Задача & Проверяемая величина & Результат проверки\\\hline
Интеграл & Ошибка на сетках $8,16,32,64$ & Совпадает с $\pi/(4n^2)$ с допуском $10^{-12}$\\\hline
Объём & Разность численного объёма и $\pi$ & Меньше $10^{-12}$ на всех сетках\\\hline
Спуск & Максимальная ошибка координат & Меньше $10^{-6}$ для обоих правил шага\\\hline
Спуск & Допустимость и изменение затрат & Все точки допустимы; роста затрат выше допуска $10^{-10}$ нет\\\hline
\end{tabular}
\end{table}
В полном отчёте команды запуска указываются для собственных скриптов.
Значения таблиц этого образца сохранены также в
\texttt{examples/report\_results.json}; рисунки находятся в
\texttt{examples/report\_figures/}.

\section*{Заключение по лабораторной работе}
В задаче 2.2 аналитически получено значение $5\pi/2$.
Численный эксперимент подтвердил уменьшение ошибки в четыре раза
при удвоении числа разбиений; объём совпал с геометрическим значением $\pi$.
В эксперименте 3.5 оба правила шага привели к точке, согласующейся
с аналитическим минимумом $(12,5+\sqrt{10})$.
Расхождения объясняются погрешностью квадратуры, критерием остановки
и машинным округлением. Эти выводы относятся только к двум рассмотренным
задачам; в полном отчёте заключение охватывает все выполненные разделы.

\section*{Источник исходных данных}
Условие лабораторной работы \texttt{main.tex}, вариант 1:
приложение А для задачи 2.2 и приложение Б для эксперимента 3.5.
В полном отчёте дополнительно указываются использованные внешние источники.
\end{document}
